\subsection{Quadratic cut grids and pointwise control on every head}
\label{actual-adaptive-pointwise}

This subsection retains the original place, normalized tower, top
lift, regular-radial identity, and $\mathrm{InDDplus2}$ hypotheses of
the compact Fourier, inverse-parameter, and block-growth proofs
in the construction record~\cite{current-construction-record}.
The coefficient is smooth and compact with a positive Gram cutoff.
All kernels $N_{j,w}$, diagonals $E_{j,w}$, quotient measures, and
positive word normalizers are the prescribed ones.  In particular,
no assertion about a discrete-series coefficient without a Gram
cutoff is included in the theorem below.

Put $s=r+1$, $N=3s$, and $M=N(2r+1)$.  Write
\[
 q_0(y)=-TT'\prod_{i=1}^k
 (y+\tfrac12+u_i)(y+\tfrac12+u'_i),\qquad
 S_w=\prod_{(\epsilon,t)\in w}\nu(t),
\]
where $\nu(t)>0$ at every permitted letter parameter.
Suppose the original cut has endpoints given by $u_\ell,u'_\ell$,
contains the interval
\[
 [-\max(u_\ell,u'_\ell)+2,
  -\min(u_\ell,u'_\ell)-2],
 \qquad |u_\ell-u'_\ell|\ge5M,
\]
and has $q_0>0$ on the nodes used below.  As always, the cut
condition includes regularity of the original operators at these
parameters.  Set $a=\min(u_\ell,u'_\ell)$ and
\[
 \mathcal Y=\{-a-3-5i:0\le i<M\}.
\]
An allowed word is an original admissible word with letters
$(+,y)$ or $(-,y+1)$, $y\in\mathcal Y$; its numerical letter
parameters are pairwise more than two apart.
Assume that
\begin{equation}\label{adaptive-normalizer-premise}
 \nu(y)\nu(y+1)\le\Lambda_0^2q_0(y)
 \quad(y\in\mathcal Y),\qquad \Lambda_0>0,
\end{equation}
and put $\Lambda=\max\{1,\Lambda_0(5/4)^{N/2}\}$.

\begin{lemma}\label{adaptive-grid-lemma}
The exact Fourier and inverse-parameter inequalities of those
proofs hold on $\mathcal Y$ with scale
$\Lambda$, for every allowed $w$ and every $K\ge0$ with
$|w|+K\le N$.  This grid fits the quadratic cut allowance
$30s^2+6$.
\end{lemma}

\begin{proof}
The smallest node is $-a-5M+2$, and the largest shifted node is
$-a-2$.  Thus every $y,y+1$ lies in the displayed cut interval.
Distinct underlying nodes are five apart.  An admissible old word
uses each underlying node at most once.  For $0\le d\le2r$,
\[
 M-|w|\ge(N-|w|)(2r+1)\ge K(d+1).
\]
Consequently one can choose $K$ disjoint blocks of $d+1$ unused
nodes.  Every choice of one node per block gives simultaneous plus
and minus extensions with all new--new and new--old numerical
distances greater than two.  For a block $y_0,\ldots,y_d$, the
leading-coefficient Lagrange weights satisfy
\[
 c_j=\prod_{i\ne j}(y_j-y_i)^{-1},\qquad
 \sum_{j=0}^d|c_j|\le(d+1)5^{-d}\le1.
\]
The last inequality includes $d=0$.

At a fresh node the exact old-word correction is
\[
 q_w(y)=q_0(y)
  \prod_{(+,t)\in w}\left(1+\frac1{y-t}\right)
  \prod_{(-,t)\in w}\left(1+\frac1{t-1-y}\right).
\]
For a negative old letter $t=z+1$, its denominator is $z-y$.
All displayed denominators therefore have modulus at least five,
and every real factor is at least $4/5$.  Hence
\[
 q_w(y)\ge(4/5)^{|w|}q_0(y)>0,\qquad
 \sqrt{\nu(y)\nu(y+1)/q_w(y)}\le\Lambda.
\]
This applies at every intermediate insertion, since its word
length is at most $N$.

The exact Capelli coefficient identities in the cited Fourier and
off-diagonal proofs hold on arbitrary injective fresh real grids.
Their respective degrees are $2j\le2r$ and
$2j-t\le2r-1$.  Apply those identities to the blocks just chosen.
The mixed determinant inequality, original wave isometry, and
identities $S_{w_+}=S_w\prod\nu(y_i)$ and
$S_{w_-}=S_w\prod\nu(y_i+1)$ give the same positive weighted
averages as in those proofs, now with the coefficient-sum bound
one and insertion ratio at most $\Lambda$.  Taking the finite
diagonal maxima proves both inequalities.  All original signs,
shifted parameters, old-word factors, and point normalizers remain
present.  Finally $5M=30s^2-15s\le30s^2+6$.
\end{proof}

\begin{theorem}\label{adaptive-all-head-pointwise}
Under these hypotheses, every valid original head $\gamma$ of
total size $s$ satisfies
\begin{equation}\label{adaptive-all-head-bound}
 E_{s,\varnothing}(\gamma)
 \le(6\Lambda^2)^s
       \max_{\substack{w\text{ allowed}\\|w|\le3s}}
                      E_{0,w}(\varnothing).
\end{equation}
More generally, appending any valid head of total size $t$ costs
at most $(6\Lambda^2)^t$ and at most $3t$ additional letters,
provided the total word budget does not exceed $N$.
\end{theorem}

\begin{proof}
Write $\mathcal W_K(w)$ for the finite set of allowed extensions
of $w$ by at most $K$ letters.  Define
$m(\Lambda,3,0)=1$ and
$m(\Lambda,3,x)=\min\{1,(\Lambda/x)^3\}$ for $x>0$.

For a singleton appended to a head $h$ of size $j$, put
$B=\max_{v\in\mathcal W_3(w)}E_{j,v}(h)$.
The Fourier inequality of Lemma~\ref{adaptive-grid-lemma}
bounds the transform of the actual continuous integrable density
$H(c)=N_{j+1,w}(h\mathbin\Vert(1,c),h\mathbin\Vert(1,c))$
by $B m(\Lambda,3,|z|)$, even when the maximizing extension
depends on $z$.  The original character is
$\exp(zc-\bar z\bar c)$ and its planar inversion constant is
$\pi^{-2}$.  Since
\[
 \int_{\mathbb C}m(\Lambda,3,|z|)\,dz=3\pi\Lambda^2,
\]
pointwise Fourier inversion for this continuous representative gives
\begin{equation}\label{adaptive-singleton-step}
 E_{j+1,w}(h\mathbin\Vert(1,c))
       \le(3/\pi)\Lambda^2B\le6\Lambda^2B.
\end{equation}
The integrable transform is established by the same explicit
majorant; no almost-everywhere replacement is used.

For a block of size $t\ge2$, let its new total level be $j$.
Suppose all $E_{j-1,v}(h\mathbin\Vert(t-1,z))$, with
$z\ne0$ and $v\in\mathcal W_3(w)$, are bounded by $B$.
The inverse-parameter inequality gives, for every fixed
$\sigma\in\mathbb C$,
\[
 \limsup_{\substack{z\to0\\z\ne0}}
 \left|N_{j-1,w}\left(h\mathbin\Vert(t-1,z),
 h\mathbin\Vert(t-1,(z^{-1}+\sigma)^{-1})\right)\right|
 \le B m(\Lambda,3,|\sigma|).
\]
The exceptional reciprocal parameter is absent from a sufficiently
small punctured neighborhood.  Thus this is the full complex
limsup required in the original block-growth theorem.

Use the primal lattice
$\mathscr L=(\pi/\Lambda)(\mathbb Z+i\mathbb Z)$.
For the prescribed character its dual is
$\Lambda(\mathbb Z+i\mathbb Z)$ and its covolume is
$\pi^2/\Lambda^2$.  Choose the translate containing the target
parameter.  Positivity permits keeping that single summand in the
block-growth inequality.  The dual sum is bounded by
\[
 B\left(1+\sum_{(a,b)\in\mathbb Z^2\setminus\{0\}}
              (a^2+b^2)^{-3/2}\right)
 \le B\left(1+8\sum_{n\ge1}n^{-2}\right)
 \le17B<6\pi B.
\]
Here the square shell $\max(|a|,|b|)=n$ has $8n$ points,
each of Euclidean norm at least $n$, and
$\sum n^{-2}\le1+\int_1^\infty x^{-2}dx=2$.
Dividing by the exact left factor
$\operatorname{covol}(\mathscr L)/\pi=\pi/\Lambda^2$
proves
\begin{equation}\label{adaptive-block-step}
 E_{j,w}(h\mathbin\Vert(t,z_0))\le6\Lambda^2B.
\end{equation}

Induct on the size of a block, with
\eqref{adaptive-singleton-step} as base and
\eqref{adaptive-block-step} as induction step.  The induction
supplies the required uniform bound on preceding diagonals, so
the finiteness premise of block growth is verified at each step.
Concatenate the blocks.  An extension of an extension is an
allowed extension, and the added lengths sum to at most $3t$.
Only finite maxima occur, so no number-of-words or number-of-cells
factor is introduced.  The empty preceding head gives
\eqref{adaptive-all-head-bound}.  For $r=0$, only the singleton
step is needed.
\end{proof}

\paragraph{Uniform scale, conditional on the spectral majorant.}
If $b\ge3$, $C\ge0$, $m=2s$, and the original spectral comparison
gives $\Lambda_0\le b^{(6+C)m}$, then
\[
 \Lambda\le b^{(7+C)m},\qquad
 (6\Lambda^2)^s\le b^{(8+C)m^2}.
\]
Indeed $(5/4)^{3s/2}\le b^{2s}$ and $6\le b^2$.
The explicit spectral comparison (NC1)--(NC15) below supplies
this premise; the adaptive grid and pointwise recursion can
therefore be applied with this scale.

\paragraph{The precise uncut endpoint.}
For an actual discrete-series coefficient $F$, let
$F_\epsilon=\chi(G/\epsilon)G^pFA$, where $\chi=0$ on
$[0,1/2]$, $\chi=1$ on $[1,\infty)$, and $A$ is the identity
or a prescribed polynomial compact-type shift.  Applying the
theorem after the established fixed-cutoff exhaustion requires
\begin{equation}\label{adaptive-uncut-endpoint}
 \limsup_{\epsilon\downarrow0}
 \max_{|w|\le3s}E_{0,w}(F_\epsilon)
       \le jN^2(G^pFA).
\end{equation}
This requires both convergence of the differentiated cutoffs and
the same-form normalized spectral word-energy inequality for
$G^pFA$.  Neither follows from the compact interpolation theorem.
For $A=1$, if this endpoint holds, orbital invariance of $G$ gives
the conditional pointwise estimate
\[
 E_F(\gamma)\le(6\Lambda^2)^sG(\gamma)^{-2p}
                   jN^2(G^pF).
\]
For equal compact types, the compact two-shift lower and contraction
upper bounds are proved in
Section~\ref{sec:archimedean-wigner-two-chain}, and their transport
to the actual ambient polynomial and radial operators is proved
in Section~\ref{actual-wigner-orbital-transport}. The endpoint
\eqref{adaptive-uncut-endpoint} follows from the spectral-word
and direct PBW cutoff theorems. Their exact hypotheses and the
final application are assembled in
Theorem~\ref{thm:archimedean-pointwise-sufficient}.

\paragraph{An explicit normalizer comparison on the original cut.}
The next estimate supplies the normalizer premise in
\eqref{adaptive-normalizer-premise}. Its exponent depends on the
rank parameter \(m=2r+2\), rather than on the number \(k\) of
compact entries. The cancellation comes from interlacing on the
two sides of the exceptional cut index.

\begin{lemma}[Two-bracket normalizer comparison]
\label{lem:archimedean-normalizer-comparison}
Let \(r\ge0\), \(k\ge1\), \(m=2r+2\), \(1\le\ell\le k\),
and let \(C,b,c\) be real numbers with
\[
 C\ge0,\qquad c\ge1,\qquad b\ge c+2,\qquad b\ge k.
\]
Let \(\lambda_1,\ldots,\lambda_{k+m-1}\) and
\(u_1,\ldots,u_k,u'_1,\ldots,u'_k\) be real numbers such that
\[
 |\lambda_j|\le b,\qquad |u_\ell|\le b,\qquad |u'_\ell|\le b,
                                                        \tag{NC1}
\]
and assume
\[
 \lambda_{j+1}+2c\le\lambda_j\qquad(1\le j\le k-2).
                                                        \tag{NC2}
\]
For each \(v=u,u'\), require the following interlacing
inequalities away from the cut index:
\[
 \begin{array}{ll}
 \lambda_j+\tfrac12\le v_j
       & (1\le j<k,\ j\ne\ell),\\
 v_{j+1}+\tfrac12\le\lambda_j
       & (1\le j<k,\ j+1\ne\ell).
 \end{array}                                            \tag{NC3}
\]
Require also the cross separation
\[
 \max(u_{j+1},u'_{j+1})+1
             \le\min(u_j,u'_j)\qquad(1\le j<k).
                                                        \tag{NC4}
\]
Suppose \(T,T'\ge b^{-Cm}\), and let \(y\) belong to the
original closed cut, meaning that either
\[
       u_\ell+2\le-y\le u'_\ell-2
\]
or the same inequalities with \(u,u'\) exchanged hold.
For a nonnegative normalizer satisfying
\[
 0\le\nu(t)\le
       \prod_{j=1}^{k+m-1}(|t+\lambda_j|+c)
       \qquad(t\in\mathbb R),
                                                        \tag{NC5}
\]
put
\[
 q_0(y)=-TT'\prod_{i=1}^k
             (y+\tfrac12+u_i)(y+\tfrac12+u'_i).
\]
Then \(q_0(y)>0\) and
\[
 \nu(y)\nu(y+1)
       \le b^{(11+2C)m}q_0(y)
       \le b^{2(6+C)m}q_0(y).                           \tag{NC6}
\]
All index ranges with an upper endpoint below their lower
endpoint are interpreted as empty.
\end{lemma}

\begin{proof}
By symmetry assume \(u_\ell+2\le-y\le u'_\ell-2\).
Set
\[
 a_j=(|y+\lambda_j|+c)(|y+1+\lambda_j|+c),\qquad
 p_i=|(y+\tfrac12+u_i)(y+\tfrac12+u'_i)|.
                                                        \tag{NC7}
\]
The cut and (NC1) give \(|y|\le b-2\), while \(b\ge3\).
Thus every \(a_j\) is nonnegative and
\[
 a_j\le(3b)^2\le b^4.                                  \tag{NC8}
\]
The cut entry has one negative and one positive factor.
For \(i<\ell\), both factors in \(p_i\) are positive,
by (NC4); for \(i>\ell\), both are negative.
Moreover each of their absolute values is at least one.
At \(i=\ell\) the two absolute values are at least
\(3/2\) and \(5/2\), respectively. In particular,
\[
 p_i\ge1\quad(1\le i\le k),\qquad
              q_0(y)=TT'\prod_{i=1}^kp_i>0.             \tag{NC9}
\]

We first record the elementary chain estimate that will be used
on each side. Suppose \(h\le b\) is a nonnegative integer and
\[
 0\le A_j\le(d_j+1+c)^2,\qquad
 d_j\ge\tfrac52+2cj,\qquad
 P_j\ge(d_j+1)^2
       \qquad(0\le j<h).
\]
Writing \(D_j=d_j+1\), the quotient \(A_j/P_j\) is at most
\((1+c/D_j)^2\). Its \(j=0\) bound is
\(K_c=(1+2c/7)^2\). For \(j\ge1\),
\[
 \left(1+\frac{c}{D_j}\right)^2
 \le\left(1+\frac1{2j}\right)^2
 \le\frac{2j+1}{2j-1}.
                                                        \tag{NC10}
\]
The last inequality is equivalent to \(4j^2-1\le4j^2\).
The odd factors telescope. For \(h\ge1\) the product is
at most \(K_c(2h-1)\), and for all \(h\ge0\) it is at most
\(K_c(2h+1)\). Since
\[
 1+\frac{2c}{7}\le\frac b2,\qquad 2h+1\le4b,
\]
we obtain the uniform chain bound
\[
                 \prod_{j=0}^{h-1}A_j
                         \le b^3\prod_{j=0}^{h-1}P_j .
                                                        \tag{NC11}
\]
The case \(h=0\) is included because \(b\ge1\).

For the left chain take the spectral indices
\(s=1,\ldots,\ell-2\) in reverse order. If this chain is
nonempty, write \(s=\ell-2-j\) and put
\[
        d_j=y+\lambda_s,\qquad A_j=a_s,\qquad P_j=p_s .
\]
The inequalities
\[
 \lambda_{\ell-2}\ge u_{\ell-1}+\tfrac12,\qquad
 u_{\ell-1}\ge u'_\ell+1,\qquad y+u'_\ell\ge2
\]
show that \(d_0\ge7/2\). Iterating (NC2) gives
\(d_j\ge7/2+2cj\), which is more than the lower bound
needed in (NC11). On the other hand, (NC3) gives
\[
 y+\tfrac12+u_s\ge d_j+1,\qquad
 y+\tfrac12+u'_s\ge d_j+1.
\]
Consequently \(p_s\ge(d_j+1)^2\), while
\[
 a_s=(d_j+c)(d_j+1+c)\le(d_j+1+c)^2.
\]
The number of indices is at most \(k\le b\), so (NC11) yields
\[
       \prod_{s=1}^{\ell-2}a_s
                          \le b^3\prod_{s=1}^{\ell-2}p_s.
                                                        \tag{NC12}
\]

For the right chain take \(s=\ell+1+j\), with
\(\ell+1\le s\le k-1\), and put
\[
 d_j=-y-\lambda_s-1,\qquad A_j=a_s,\qquad P_j=p_{s+1}.
\]
If the chain is nonempty, (NC3), (NC4), and the cut give
\[
 \lambda_{\ell+1}\le u_{\ell+1}-\tfrac12
                         \le u_\ell-\tfrac32
                         \le-y-\tfrac72.
\]
Hence \(d_0\ge5/2\), and (NC2) gives \(d_j\ge5/2+2cj\).
Also (NC3) gives
\[
 y+\tfrac12+u_{s+1}\le-(d_j+1),\qquad
 y+\tfrac12+u'_{s+1}\le-(d_j+1).
\]
Thus the same chain estimate applies and gives
\[
       \prod_{s=\ell+1}^{k-1}a_s
                          \le b^3\prod_{i=\ell+2}^{k}p_i.
                                                        \tag{NC13}
\]

The two spectral chains are disjoint. Among the first \(k-1\)
spectral indices they leave at most the two indices
\(\ell-1,\ell\); there are \(m\) further spectral indices
from \(k\) through \(k+m-1\). Thus at most \(m+2=2r+4\)
indices remain unpaired. Apply (NC8) to these indices and
(NC12)--(NC13) to the chains. By (NC9), inserting any missing
compact factors only increases the product. We obtain
\[
 \begin{split}
 \prod_{j=1}^{k+m-1}a_j
   &\le b^{4(m+2)+6}\prod_{i=1}^kp_i\\
   &=b^{8r+22}\prod_{i=1}^kp_i
     \le b^{11m}\prod_{i=1}^kp_i .
 \end{split}                                           \tag{NC14}
\]
The last inequality uses \(8r+22\le11(2r+2)\).
For \(k=1\), both chains are empty and there are only \(m\)
spectral indices; the same upper bound applies. For
\(\ell=1\) or \(\ell=k\), one chain is empty and at most
one exceptional index lies among the first \(k-1\).
Thus no boundary index has been used outside its stated range.

Finally (NC5), \(TT'\ge b^{-2Cm}>0\), and (NC9) imply
\[
 \nu(y)\nu(y+1)
 \le\prod_{j=1}^{k+m-1}a_j
 \le b^{11m}(TT')^{-1}q_0(y)
 \le b^{(11+2C)m}q_0(y).
\]
Since \(b\ge1\) and \(m\ge2\), this implies the stated weaker
exponent \(2(6+C)m\) as well.
\end{proof}

\begin{corollary}[Explicit scale for the quadratic cut grid]
\label{cor:adaptive-explicit-normalizer}
Under the hypotheses of the preceding lemma at every node of
the original quadratic cut grid, the premise
\eqref{adaptive-normalizer-premise} holds with
\[
                    \Lambda_0=b^{(6+C)m}.
\]
Retain the tower, top lift, regular-radial identity and
\(\mathrm{InDDplus2}\) hypotheses of
Theorem~\ref{adaptive-all-head-pointwise}. With \(s=r+1\),
\(m=2s\), and word budget \(N=3s\), its conclusion becomes
\[
 E_{s,\varnothing}(\gamma)
      \le b^{(8+C)m^2}
           \max_{\substack{w\ {\rm allowed}\\|w|\le3s}}
                           E_{0,w}(\varnothing).
                                                        \tag{NC15}
\]
\end{corollary}
\begin{proof}
The first assertion is (NC6). The scale in that theorem is
\(\Lambda=b^{(6+C)m}(5/4)^{3s/2}\), since its factors exceed one.
Its cost is therefore
\[
 (6\Lambda^2)^s
    =6^s b^{4(6+C)s^2}(5/4)^{3s^2}
    \le b^{(27+4C)s^2+2s}
    \le b^{(8+C)m^2}.
\]
Here \(b\ge3\) gives \(6\le b^2\) and \(5/4\le b\), and
\(s\ge1\) gives \(2s\le5s^2\). No uncut coefficient estimate
or removal of the Gram cutoff is implicit in this corollary.
\end{proof}
